% ---------------------------------------------------------------------------
% Préambule commun du manuscrit reactant-analyzer.
% Compilé avec pdflatex. Les paquets non essentiels sont chargés sous garde
% \IfFileExists pour que le document compile aussi sur une installation TeX
% minimale (mode dégradé : pas de TikZ, pas de boîtes colorées).
% ---------------------------------------------------------------------------
\usepackage[utf8]{inputenc}
\usepackage[T1]{fontenc}
\usepackage{lmodern}
\IfFileExists{zi4.sty}{\usepackage[scaled=0.9]{zi4}}{}          % Inconsolata pour le code si dispo
\usepackage{microtype}
\IfFileExists{french.ldf}
  {\usepackage[french]{babel}\frenchsetup{StandardLists=true}}
  {\usepackage[english]{babel}\providecommand{\og}{\guillemotleft~}\providecommand{\fg}{~\guillemotright}}
\usepackage[a4paper,margin=2.6cm,top=2.8cm,bottom=2.8cm,headheight=15pt]{geometry}
\usepackage{amsmath,amssymb,amsthm,mathtools}
\usepackage{xcolor}
\usepackage{graphicx}
\usepackage{array,booktabs,longtable,tabularx}
\IfFileExists{multirow.sty}{\usepackage{multirow}}{}
\IfFileExists{enumitem.sty}{\usepackage{enumitem}\setlist{noitemsep,topsep=3pt}}{}
\usepackage{caption}
\usepackage{subcaption}
\usepackage{float}
\usepackage{etoolbox}
\IfFileExists{needspace.sty}{\usepackage{needspace}}{\providecommand{\needspace}[1]{}}
\usepackage{fancyhdr}
\usepackage{listings}
\IfFileExists{algorithm2e.sty}{\usepackage[french,algoruled,linesnumbered,vlined]{algorithm2e}}{}
\IfFileExists{imakeidx.sty}{\usepackage{imakeidx}}{\usepackage{makeidx}}
\makeindex
\IfFileExists{csquotes.sty}{\usepackage{csquotes}}{}
\usepackage[hidelinks,pdfusetitle,bookmarksnumbered,bookmarksopen,bookmarksopenlevel=1]{hyperref}
\IfFileExists{cleveref.sty}{\usepackage[french,nameinlink,noabbrev]{cleveref}}{}
% Le « : » actif de babel-french casse \cref sur les labels chap:/sec:/def:.
% On le désactive pour tout le document ; le rendu « mot : suite » reste
% correct avec une espace ordinaire. Les chapitres n'ont donc pas à le faire.
\AtBeginDocument{\shorthandoff{:}}
\AtBeginDocument{\renewcommand{\shorthandon}[1]{}}
% Noms cleveref des boîtes tcolorbox (sinon « ?? »)
\makeatletter
  \@ifpackageloaded{cleveref}{%
    \crefname{tcb@cnt@definition}{définition}{définitions}%
    \Crefname{tcb@cnt@definition}{Définition}{Définitions}%
    \crefname{tcb@cnt@theoreme}{théorème}{théorèmes}%
    \Crefname{tcb@cnt@theoreme}{Théorème}{Théorèmes}%
    \crefname{tcb@cnt@proposition}{proposition}{propositions}%
    \Crefname{tcb@cnt@proposition}{Proposition}{Propositions}%
    \crefname{tcb@cnt@lemme}{lemme}{lemmes}%
    \Crefname{tcb@cnt@lemme}{Lemme}{Lemmes}%
    \crefname{tcb@cnt@invariant}{invariant}{invariants}%
    \Crefname{tcb@cnt@invariant}{Invariant}{Invariants}%
    \crefname{tcb@cnt@exemple}{exemple}{exemples}%
    \Crefname{tcb@cnt@exemple}{Exemple}{Exemples}%
    \crefname{tcb@cnt@exercice}{exercice}{exercices}%
    \Crefname{tcb@cnt@exercice}{Exercice}{Exercices}%
    \crefname{lstlisting}{extrait}{extraits}%
    \Crefname{lstlisting}{Extrait}{Extraits}%
    \crefname{algocf}{algorithme}{algorithmes}%
    \Crefname{algocf}{Algorithme}{Algorithmes}%
  }{}%
\makeatother

% ---- TikZ (schémas) --------------------------------------------------------
\IfFileExists{tikz.sty}{
  \usepackage{tikz}
  \usetikzlibrary{arrows.meta,positioning,shapes.geometric,shapes.misc,calc,fit,matrix,chains,decorations.pathreplacing,backgrounds,automata}
  \newif\ifhastikz\hastikztrue
}{
  \newif\ifhastikz\hastikzfalse
}

% ---- Couleurs du manuscrit -------------------------------------------------
\definecolor{ra-blue}{HTML}{1F4E79}
\definecolor{ra-teal}{HTML}{0F766E}
\definecolor{ra-amber}{HTML}{B45309}
\definecolor{ra-red}{HTML}{9F1239}
\definecolor{ra-gray}{HTML}{4B5563}
\definecolor{ra-bg}{HTML}{F7F7F5}
\definecolor{ra-kw}{HTML}{7C3AED}
\definecolor{ra-str}{HTML}{047857}
\definecolor{ra-cmt}{HTML}{6B7280}
\definecolor{ra-num}{HTML}{B45309}

% ---- Boîtes (tcolorbox si dispo, sinon amsthm) -----------------------------
\theoremstyle{definition}
\newif\ifhastcb
\IfFileExists{tcolorbox.sty}{\hastcbtrue}{\hastcbfalse}
\ifhastcb
  \usepackage[most]{tcolorbox}
  \tcbset{
    manuscritbox/.style={enhanced,breakable,colback=ra-bg,colframe=#1,
      boxrule=0.6pt,arc=1.5pt,left=6pt,right=6pt,top=4pt,bottom=4pt,
      fonttitle=\bfseries\small,coltitle=white,colbacktitle=#1,
      attach boxed title to top left={yshift=-2mm,xshift=4mm},
      boxed title style={arc=1pt,boxrule=0pt},before skip=8pt,after skip=8pt},
  }
  \newtcbtheorem[number within=chapter]{definition}{Définition}{manuscritbox=ra-blue}{def}
  \newtcbtheorem[number within=chapter]{theoreme}{Théorème}{manuscritbox=ra-blue}{thm}
  \newtcbtheorem[number within=chapter]{proposition}{Proposition}{manuscritbox=ra-blue}{prop}
  \newtcbtheorem[number within=chapter]{lemme}{Lemme}{manuscritbox=ra-blue}{lem}
  \newtcbtheorem[number within=chapter]{invariant}{Invariant}{manuscritbox=ra-teal}{inv}
  \newtcbtheorem[number within=chapter]{exemple}{Exemple}{manuscritbox=ra-teal,colback=white}{ex}
  \newtcbtheorem[number within=chapter]{exercice}{Exercice}{manuscritbox=ra-gray,colback=white}{exo}
  \newtcolorbox{remarque}{manuscritbox=ra-gray,title=Remarque}
  \newtcolorbox{piege}{manuscritbox=ra-red,title=Piège}
  \newtcolorbox{react}[1][]{manuscritbox=ra-amber,title={Rappel React\ifstrempty{#1}{}{ : #1}}}
  \newtcolorbox{decision}[1][]{manuscritbox=ra-blue,colback=white,title={Décision\ifstrempty{#1}{}{ : #1}}}
  \newtcolorbox{solution}{manuscritbox=ra-gray,colback=white,title=Solution}
\else
  \newtheorem{definitionthm}{Définition}[chapter]
  \newtheorem{theoremethm}[definitionthm]{Théorème}
  \newtheorem{propositionthm}[definitionthm]{Proposition}
  \newtheorem{lemmethm}[definitionthm]{Lemme}
  \newtheorem{invariantthm}[definitionthm]{Invariant}
  \newtheorem{exemplethm}[definitionthm]{Exemple}
  \newtheorem{exercicethm}[definitionthm]{Exercice}
  % Même interface que tcolorbox : \begin{definition}{Titre}{label}
  \newenvironment{definition}[2]{\begin{definitionthm}[#1]\label{def:#2}}{\end{definitionthm}}
  \newenvironment{theoreme}[2]{\begin{theoremethm}[#1]\label{thm:#2}}{\end{theoremethm}}
  \newenvironment{proposition}[2]{\begin{propositionthm}[#1]\label{prop:#2}}{\end{propositionthm}}
  \newenvironment{lemme}[2]{\begin{lemmethm}[#1]\label{lem:#2}}{\end{lemmethm}}
  \newenvironment{invariant}[2]{\begin{invariantthm}[#1]\label{inv:#2}}{\end{invariantthm}}
  \newenvironment{exemple}[2]{\begin{exemplethm}[#1]\label{ex:#2}}{\end{exemplethm}}
  \newenvironment{exercice}[2]{\begin{exercicethm}[#1]\label{exo:#2}}{\end{exercicethm}}
  \newenvironment{remarque}{\par\medskip\noindent\textbf{Remarque.}\ }{\par\medskip}
  \newenvironment{piege}{\par\medskip\noindent\textbf{Piège.}\ }{\par\medskip}
  \newenvironment{react}[1][]{\par\medskip\noindent\textbf{Rappel React\ifstrempty{#1}{}{ : #1}.}\ }{\par\medskip}
  \newenvironment{decision}[1][]{\par\medskip\noindent\textbf{Décision\ifstrempty{#1}{}{ : #1}.}\ }{\par\medskip}
  \newenvironment{solution}{\par\medskip\noindent\textbf{Solution.}\ }{\par\medskip}
\fi

% ---- En-têtes --------------------------------------------------------------
\pagestyle{fancy}
\fancyhf{}
\fancyhead[LE,RO]{\small\thepage}
\fancyhead[RE]{\small\nouppercase{\leftmark}}
\fancyhead[LO]{\small\nouppercase{\rightmark}}
\renewcommand{\headrulewidth}{0.3pt}
\fancypagestyle{plain}{\fancyhf{}\fancyfoot[C]{\small\thepage}\renewcommand{\headrulewidth}{0pt}}

\setcounter{tocdepth}{2}
\setcounter{secnumdepth}{3}
\setlength{\parskip}{3pt plus 1pt}
\raggedbottom
\emergencystretch=2em

% ---------------------------------------------------------------------------
% Configuration de listings : Rust, TypeScript/TSX, JSON, shell.
% Les extraits de code sont pris VERBATIM dans le dépôt (\rustsnippet) ou
% inclus inline. Le code source contient des caractères Unicode (⊤, →, ─,
% é, …) : la table `literate` ci-dessous les rend sous pdflatex.
% ---------------------------------------------------------------------------

% Caractères Unicode utilisables dans la PROSE (hors listings)
\DeclareUnicodeCharacter{22A4}{\ensuremath{\top}}
\DeclareUnicodeCharacter{22A5}{\ensuremath{\bot}}
\DeclareUnicodeCharacter{2293}{\ensuremath{\sqcap}}
\DeclareUnicodeCharacter{2294}{\ensuremath{\sqcup}}
\DeclareUnicodeCharacter{2291}{\ensuremath{\sqsubseteq}}
\DeclareUnicodeCharacter{2292}{\ensuremath{\sqsupseteq}}
\DeclareUnicodeCharacter{22E2}{\ensuremath{\not\sqsubseteq}}
\DeclareUnicodeCharacter{2192}{\ensuremath{\rightarrow}}
\DeclareUnicodeCharacter{2190}{\ensuremath{\leftarrow}}
\DeclareUnicodeCharacter{2194}{\ensuremath{\leftrightarrow}}
\DeclareUnicodeCharacter{21D2}{\ensuremath{\Rightarrow}}
\DeclareUnicodeCharacter{27F9}{\ensuremath{\Longrightarrow}}
\DeclareUnicodeCharacter{27FA}{\ensuremath{\Longleftrightarrow}}
\DeclareUnicodeCharacter{21C4}{\ensuremath{\rightleftarrows}}
\DeclareUnicodeCharacter{2264}{\ensuremath{\leq}}
\DeclareUnicodeCharacter{2265}{\ensuremath{\geq}}
\DeclareUnicodeCharacter{2260}{\ensuremath{\neq}}
\DeclareUnicodeCharacter{2261}{\ensuremath{\equiv}}
\DeclareUnicodeCharacter{221E}{\ensuremath{\infty}}
\DeclareUnicodeCharacter{2200}{\ensuremath{\forall}}
\DeclareUnicodeCharacter{2203}{\ensuremath{\exists}}
\DeclareUnicodeCharacter{2227}{\ensuremath{\wedge}}
\DeclareUnicodeCharacter{2228}{\ensuremath{\vee}}
\DeclareUnicodeCharacter{222A}{\ensuremath{\cup}}
\DeclareUnicodeCharacter{2229}{\ensuremath{\cap}}
\DeclareUnicodeCharacter{2208}{\ensuremath{\in}}
\DeclareUnicodeCharacter{2209}{\ensuremath{\notin}}
\DeclareUnicodeCharacter{2286}{\ensuremath{\subseteq}}
\DeclareUnicodeCharacter{2282}{\ensuremath{\subset}}
\DeclareUnicodeCharacter{2205}{\ensuremath{\emptyset}}
\DeclareUnicodeCharacter{2295}{\ensuremath{\oplus}}
\DeclareUnicodeCharacter{2213}{\ensuremath{\mp}}
\DeclareUnicodeCharacter{00B1}{\ensuremath{\pm}}
\DeclareUnicodeCharacter{00D7}{\ensuremath{\times}}
\DeclareUnicodeCharacter{00B7}{\ensuremath{\cdot}}
\DeclareUnicodeCharacter{2022}{\textbullet}
\DeclareUnicodeCharacter{2026}{...}
\DeclareUnicodeCharacter{2500}{--}
\DeclareUnicodeCharacter{2713}{\ensuremath{\checkmark}}
\DeclareUnicodeCharacter{2705}{\ensuremath{\checkmark}}
\DeclareUnicodeCharacter{274C}{\ensuremath{\times}}
\DeclareUnicodeCharacter{26A0}{\textbf{!}}
\DeclareUnicodeCharacter{2139}{\textbf{i}}
\DeclareUnicodeCharacter{03B3}{\ensuremath{\gamma}}
\DeclareUnicodeCharacter{03B1}{\ensuremath{\alpha}}
\DeclareUnicodeCharacter{2081}{\ensuremath{{}_1}}
\DeclareUnicodeCharacter{1D62}{\ensuremath{{}_i}}
\DeclareUnicodeCharacter{208B}{\ensuremath{{}_-}}
\DeclareUnicodeCharacter{2207}{\ensuremath{\nabla}}

\lstdefinelanguage{Rust}{
  keywords={as,async,await,break,const,continue,crate,dyn,else,enum,extern,
    false,fn,for,if,impl,in,let,loop,match,mod,move,mut,pub,ref,return,
    self,Self,static,struct,super,trait,true,type,unsafe,use,where,while,
    macro_rules,union},
  ndkeywords={u8,u16,u32,u64,u128,usize,i8,i16,i32,i64,i128,isize,f32,f64,
    bool,char,str,String,Vec,Option,Some,None,Result,Ok,Err,Box,Rc,Arc,
    RefCell,Cell,HashMap,HashSet,BTreeMap,BTreeSet,VecDeque,FxHashMap,
    IndexMap,IndexSet,Cow,PhantomData},
  sensitive=true,
  comment=[l]{//},
  morecomment=[s]{/*}{*/},
  morestring=[b]",
  alsoletter={!},
  morekeywords=[3]{debug_assert!,assert!,assert_eq!,assert_ne!,panic!,
    unreachable!,todo!,vec!,format!,write!,writeln!,println!,eprintln!,
    matches!,include_str!},
  keywordstyle=[3]\color{ra-teal},
}

\lstdefinelanguage{TypeScript}{
  keywords={const,let,var,function,return,if,else,for,while,do,break,continue,
    new,delete,typeof,instanceof,in,of,switch,case,default,throw,try,catch,
    finally,import,export,from,as,class,extends,implements,interface,type,
    enum,async,await,yield,this,super,null,undefined,true,false,void,never,
    any,unknown,number,string,boolean,readonly,public,private,protected,
    static,declare,namespace,module,keyof,satisfies},
  ndkeywords={useState,useEffect,useMemo,useCallback,useRef,useContext,
    useReducer,useLayoutEffect,useId,useTransition,useSyncExternalStore,
    useDeferredValue,useImperativeHandle,memo,forwardRef,createContext,lazy,
    Suspense,React},
  sensitive=true,
  comment=[l]{//},
  morecomment=[s]{/*}{*/},
  morestring=[b]',
  morestring=[b]",
  morestring=[b]`,
}
\lstalias{TSX}{TypeScript}
\lstalias{JavaScript}{TypeScript}
\lstalias{JSX}{TypeScript}

\lstdefinelanguage{JSON}{
  keywords={true,false,null},
  sensitive=true,
  morestring=[b]",
  comment=[l]{//},
}

\lstdefinestyle{manuscrit}{
  basicstyle=\ttfamily\footnotesize,
  keywordstyle=\color{ra-kw},
  ndkeywordstyle=\color{ra-blue},
  stringstyle=\color{ra-str},
  commentstyle=\color{ra-cmt}\itshape,
  numberstyle=\tiny\color{ra-gray},
  numbers=left,
  numbersep=7pt,
  stepnumber=1,
  showstringspaces=false,
  breaklines=true,
  breakatwhitespace=false,
  postbreak=\mbox{\textcolor{ra-gray}{$\hookrightarrow$}\space},
  tabsize=4,
  columns=fullflexible,
  keepspaces=true,
  frame=single,
  framerule=0.3pt,
  rulecolor=\color{ra-gray!50},
  backgroundcolor=\color{ra-bg},
  xleftmargin=1.6em,
  framexleftmargin=1.4em,
  captionpos=t,
  abovecaptionskip=2pt,
  belowcaptionskip=2pt,
  aboveskip=8pt,
  belowskip=8pt,
  extendedchars=true,
  inputencoding=utf8,
  upquote=false,
  literate=
    {─}{{-}}1 {━}{{=}}1 {│}{{|}}1 {┌}{{+}}1 {┐}{{+}}1 {└}{{+}}1 {┘}{{+}}1
    {├}{{+}}1 {┤}{{+}}1 {┬}{{+}}1 {┴}{{+}}1 {┼}{{+}}1 {▶}{{>}}1 {▸}{{>}}1
    {—}{{---}}1 {–}{{--}}1 {…}{{...}}3 {•}{{\textbullet}}1 {·}{{$\cdot$}}1
    {→}{{$\rightarrow$}}1 {←}{{$\leftarrow$}}1 {↔}{{$\leftrightarrow$}}1
    {⇒}{{$\Rightarrow$}}1 {⟹}{{$\Longrightarrow$}}2 {⟺}{{$\Longleftrightarrow$}}2
    {⇄}{{$\rightleftarrows$}}1 {↦}{{$\mapsto$}}1
    {⊤}{{$\top$}}1 {⊥}{{$\bot$}}1 {⊓}{{$\sqcap$}}1 {⊔}{{$\sqcup$}}1
    {⊑}{{$\sqsubseteq$}}1 {⊒}{{$\sqsupseteq$}}1 {⋢}{{$\not\sqsubseteq$}}1
    {∇}{{$\nabla$}}1 {∞}{{$\infty$}}1 {∀}{{$\forall$}}1 {∃}{{$\exists$}}1
    {∧}{{$\wedge$}}1 {∨}{{$\vee$}}1 {∪}{{$\cup$}}1 {∩}{{$\cap$}}1
    {∈}{{$\in$}}1 {∉}{{$\notin$}}1 {⊆}{{$\subseteq$}}1 {⊂}{{$\subset$}}1
    {∅}{{$\emptyset$}}1 {⊕}{{$\oplus$}}1 {∓}{{$\mp$}}1 {±}{{$\pm$}}1
    {×}{{$\times$}}1 {≤}{{$\leq$}}1 {≥}{{$\geq$}}1 {≠}{{$\neq$}}1 {≡}{{$\equiv$}}1
    {γ}{{$\gamma$}}1 {α}{{$\alpha$}}1 {β}{{$\beta$}}1 {λ}{{$\lambda$}}1
    {ρ}{{$\rho$}}1 {σ}{{$\sigma$}}1 {ᵢ}{{$_i$}}1 {₁}{{$_1$}}1 {₋}{{$_-$}}1
    {✓}{{$\checkmark$}}1 {✅}{{$\checkmark$}}1 {❌}{{$\times$}}1 {⚠}{{!}}1 {ℹ}{{i}}1
    {§}{{\S}}1 {°}{{$^\circ$}}1 {’}{{'}}1 {‘}{{`}}1 {“}{{``}}1 {”}{{''}}1
    {«}{{\guillemotleft}}1 {»}{{\guillemotright}}1
    {é}{{\'e}}1 {è}{{\`e}}1 {ê}{{\^e}}1 {ë}{{\"e}}1 {à}{{\`a}}1 {â}{{\^a}}1
    {î}{{\^i}}1 {ï}{{\"i}}1 {ô}{{\^o}}1 {ö}{{\"o}}1 {û}{{\^u}}1 {ù}{{\`u}}1
    {ü}{{\"u}}1 {ç}{{\c{c}}}1 {ñ}{{\~n}}1 {É}{{\'E}}1 {È}{{\`E}}1 {À}{{\`A}}1
    {Ç}{{\c{C}}}1 {œ}{{\oe}}1 {æ}{{\ae}}1,
}
\lstset{style=manuscrit,language=Rust}
\renewcommand{\lstlistingname}{Extrait}
\renewcommand{\lstlistlistingname}{Liste des extraits de code}

% ---------------------------------------------------------------------------
% Macros du manuscrit. Les chapitres n'en définissent AUCUNE autre.
% ---------------------------------------------------------------------------

% Racine du dépôt, relative à docs/manuscrit (pour \rustsnippet & co).
\newcommand{\repoRoot}{../..}
% Commit photographié par le manuscrit.
\newcommand{\snapshotCommit}{e67b10a}
\newcommand{\snapshotDate}{27 septembre 2026}

% ---- Texte -----------------------------------------------------------------
% Identifiant / chemin en machine à écrire, accepte _ # & % $ ^ ~ sans
% échappement ET les caractères UTF-8 (⊤, ✓, é…). Principe : on détokenise
% l'argument, puis on le relit (\scantokens) avec les caractères spéciaux
% rendus inertes, tandis que les octets UTF-8 restent actifs et sont décodés
% par inputenc. \{ \} \# \ldots restent des commandes et s'impriment.
\makeatletter
\newcommand{\ra@codecatcodes}{%
  \catcode`\_=12 \catcode`\#=12 \catcode`\&=12 \catcode`\%=12
  \catcode`\$=12 \catcode`\^=12 \catcode`\~=12 \endlinechar=-1 }
\DeclareRobustCommand{\code}[1]{\begingroup\ra@codecatcodes\edef\ra@tmp{\detokenize{#1}}\texttt{\scantokens\expandafter{\ra@tmp}}\endgroup}
\DeclareRobustCommand{\file}[1]{\begingroup\ra@codecatcodes\edef\ra@tmp{\detokenize{#1}}\texttt{\scantokens\expandafter{\ra@tmp}}\endgroup}
% Dans les signets PDF (hyperref), \code et \file rendent leur argument tel quel.
\pdfstringdefDisableCommands{\let\code\@firstofone \let\file\@firstofone}
\makeatother
% Référence à un endroit du code : \src{src/ir/expr.rs}{12}{40}  /  \srcline{src/ir/expr.rs}{12}
\newcommand{\src}[3]{\file{#1}, l.~#2--#3}
\newcommand{\srcline}[2]{\file{#1}, l.~#2}
% Nom d'une règle du catalogue : \regle{infinite-loop}
\newcommand{\regle}[1]{\textsf{\detokenize{#1}}}
% Nom d'une relation du moteur : \relation{writers}
\newcommand{\relation}[1]{\textsf{\detokenize{#1}}}
% Référence à une ADR : \adr{27} -> lien vers l'annexe, si le label existe
\newcommand{\adr}[1]{\hyperref[adr:#1]{ADR-#1}}
% Issue GitHub : \issue{63}
\newcommand{\issue}[1]{\href{https://github.com/rboudrouss/reactant-analyzer/issues/#1}{\##1}}
% Terme défini pour la première fois (italique + index)
\newcommand{\terme}[1]{\emph{#1}\index{#1}}
% Entrée d'index sans mise en forme
\newcommand{\idx}[1]{\index{#1}}
% Niveaux de diagnostic
\newcommand{\Error}{\textsc{Error}}
\newcommand{\Warning}{\textsc{Warning}}
\newcommand{\Info}{\textsc{Info}}
% Nom du projet
\newcommand{\reactant}{\textsf{reactant-analyzer}}
\newcommand{\React}{React}

% ---- Mathématiques ---------------------------------------------------------
\newcommand{\must}{\mathsf{must}}
\newcommand{\may}{\mathsf{may}}
\newcommand{\Top}{\top}
\newcommand{\Bot}{\bot}
\newcommand{\join}{\sqcup}
\newcommand{\meet}{\sqcap}
\newcommand{\widen}{\mathbin{\nabla}}
\newcommand{\lleq}{\sqsubseteq}
\newcommand{\lgeq}{\sqsupseteq}
\newcommand{\lfp}{\mathrm{lfp}}
\newcommand{\gam}{\gamma}
\newcommand{\abs}{\alpha}
\newcommand{\sem}[1]{\llbracket #1 \rrbracket}
\providecommand{\llbracket}{[\![}
\providecommand{\rrbracket}{]\!]}
\newcommand{\Int}{\mathbb{Z}}
\newcommand{\Nat}{\mathbb{N}}
\newcommand{\Bool}{\mathbb{B}}
\newcommand{\powerset}{\mathcal{P}}
\newcommand{\Stab}{\mathcal{S}}
\newcommand{\Val}{\mathcal{V}}
\newcommand{\Env}{\mathcal{E}}
\newcommand{\Heap}{\mathcal{H}}
\newcommand{\Store}{\mathcal{M}}
\newcommand{\Render}{\mathsf{render}}
\newcommand{\Commit}{\mathsf{commit}}
\newcommand{\Effect}{\mathsf{effect}}
\newcommand{\Tick}{\mathsf{tick}}
\newcommand{\Fix}{\mathsf{fix}}
\DeclareMathOperator{\dom}{dom}
\DeclareMathOperator{\reach}{reach}
\DeclareMathOperator{\pred}{pred}
\DeclareMathOperator{\succs}{succ}

% ---- Extraits de code pris dans le dépôt -----------------------------------
% \rustsnippet[options listings]{chemin relatif au dépôt}{première ligne}{dernière ligne}{légende}
% Inclut VERBATIM les lignes du fichier, numérotées avec leurs vrais numéros.
\newcommand{\rustsnippet}[5][]{%
  \lstinputlisting[language=Rust,firstline=#3,lastline=#4,firstnumber=#3,
    caption={\file{#2}, l.~#3--#4 --- #5},#1]{\repoRoot/#2}}
% Même chose pour un fichier TSX/TS/JS du dépôt (fixtures, tests)
\newcommand{\tsxsnippet}[5][]{%
  \lstinputlisting[language=TypeScript,firstline=#3,lastline=#4,firstnumber=#3,
    caption={\file{#2}, l.~#3--#4 --- #5},#1]{\repoRoot/#2}}
% Fichier entier
\newcommand{\rustfile}[3][]{\lstinputlisting[language=Rust,caption={\file{#2} --- #3},#1]{\repoRoot/#2}}
\newcommand{\tsxfile}[3][]{\lstinputlisting[language=TypeScript,caption={\file{#2} --- #3},#1]{\repoRoot/#2}}
% Sortie d'une commande (texte brut, sans numéros) : \begin{sortie} ... \end{sortie}
\lstnewenvironment{sortie}[1][]{\lstset{language={},numbers=none,basicstyle=\ttfamily\scriptsize,backgroundcolor=\color{white},frame=leftline,#1}}{}
% Terminal
\lstnewenvironment{shell}[1][]{\lstset{language=bash,numbers=none,basicstyle=\ttfamily\footnotesize,#1}}{}

% ---- Structure de chapitre -------------------------------------------------
% Encadré d'objectifs en tête de chapitre
\newenvironment{objectifs}{%
  \par\noindent\begin{minipage}{\linewidth}\small
  \textbf{Objectifs du chapitre.}\begin{itemize}}{\end{itemize}\end{minipage}\par\bigskip}
% Résumé de fin de chapitre
\newenvironment{resume}{\section*{Résumé du chapitre}\addcontentsline{toc}{section}{Résumé du chapitre}}{}
% Prérequis
\newcommand{\prerequis}[1]{\par\noindent{\small\textbf{Prérequis :} #1}\par\medskip}

% ---- Figures ---------------------------------------------------------------
% Style TikZ commun (si TikZ est chargé)
\ifhastikz
\tikzset{
  bloc/.style={draw=ra-blue,fill=ra-bg,rounded corners=2pt,thick,align=center,inner sep=5pt,font=\small},
  noeud/.style={draw=ra-blue,circle,thick,inner sep=2pt,font=\small},
  cfgbloc/.style={draw=ra-gray,fill=white,rectangle,align=left,inner sep=4pt,font=\ttfamily\scriptsize},
  fleche/.style={-{Stealth[length=2mm]},thick,draw=ra-gray},
  flechemust/.style={fleche,draw=ra-blue},
  flechemay/.style={fleche,dashed,draw=ra-amber},
  treillis/.style={font=\small},
  etiquette/.style={font=\scriptsize\itshape,text=ra-gray},
}
\fi

